% TOP_PROBLEM: {"id": 2307043, "title": "Research Problems in Function Theory — Problem 7.43", "category_id": 8, "category": {"id": 8, "name": "analysis", "display_name": "Analysis", "description": "Limits, continuity, calculus, and function theory.", "slug": "analysis", "order_index": 8, "created_at": "2026-07-31T15:26:25.671Z"}, "status": "open"}
% FIRST_POSTED: null
% ENTRY_KIND: statement_only
% Imported from: batch05.tex; UnsolvedMath ID 2307043
\documentclass[11pt]{article}
\usepackage{amsmath,amssymb,mathtools}
\usepackage{fontspec,unicode-math}
\setmainfont{Latin Modern Roman}
\setmathfont{Latin Modern Math}
\usepackage{xurl,hyperref}
\newcommand{\sref}[2]{\hyperlink{source:#1:#2}{[S#2]}}
\newcommand{\cataloguescope}[1]{\paragraph{Mathematical statement.}#1}
\begin{document}
% UnsolvedMath ID 2307043; source catalogue code AMR-022-7043
\setcounter{section}{2307042}
\section{Research Problems in Function Theory — Problem 7.43}\label{2307043}
{\small\sffamily UnsolvedMath ID 2307043 · Analysis}
\subsection{Definitions and mathematical statement}

\paragraph{Shared notation.}$\mathbb N=\{1,2,\ldots\}$, and $\mathbb Z,\mathbb Q,\mathbb R,\mathbb C$ denote the integers, rationals, reals and complex numbers. Any other conventions are specified locally.


Let $D\subset\mathbb C$ be a bounded simply connected domain. A conformal self-map means a biholomorphic map $f:D\to D$. Call $D$ conformally rigid if there is $\varepsilon>0$ such that every such $f$ satisfying $|f(z)-z|<\varepsilon$ for every $z\in D$ is the identity. For a prime end represented by a nested chain of crosscuts, its impression is the intersection of the closures of the associated nested end domains. Saying the prime end is a singleton means that this impression is a single boundary point.

\paragraph{Statement recorded in the supplied source.}
Must every bounded simply connected domain that is not conformally rigid have singleton impressions for all its prime ends? Equivalently, is non-rigidity sufficient for every prime end to be a singleton, as well as being necessary under the source's stated implication? \sref{2307043}{2}
The printed premise says singleton prime ends imply non-rigidity; the requested converse is the reverse implication. Update7.43 records Gaier's counterexample to the actual converse. \sref{2307043}{2}
\subsection{Short English statement}
Does failure of conformal rigidity force every prime end to have a single-point impression?
\subsection{Sources}
\begin{description}
\item[\hypertarget{source:2307043:1}{S1}] ULAM.AI and the UnsolvedMath Contributors, \emph{UnsolvedMath: A Curated Collection of Open Mathematics Problems}, record 2307043 (AMR-022-7043). CC BY4.0. \url{https://huggingface.co/datasets/ulamai/UnsolvedMath}.
\item[\hypertarget{source:2307043:2}{S2}] W. K. Hayman and E. F. Lingham, \emph{Research Problems in Function Theory} (2018), Chapter7, Problem7.43 and its complete update. \url{https://arxiv.org/abs/1809.07200}.
\end{description}
\label{end:2307043}
\end{document}
